% ======================================================================
% Ontologie der Kontinua -- Core 1.3
% Cross-K-Modul: crossk_k5_k6.tex
% Cross-Level-Relationen zwischen K5 und K6
% ======================================================================

\subsubsection{K5–K6-Querverbindung}
\label{sec:crossk-k5-k6}

Der Übergang $K_5 \rightarrow K_6$ markiert die Entstehung des kognitiven
Kontinuums aus dem neuronalen Kontinuum. Während $K_5$ elektrische Erregbarkeit,
Spike-Propagation und netzwerkartige Synchronisation unterstützt, führt $K_6$ ein:
\begin{itemize}
  \item konzeptionelle Achsen und semantische Unterscheidungen,
  \item kognitive Potentiale (Sicherheit, Relevanz, Salienz),
  \item Schwellen für Kohärenz und Widerspruch,
  \item strukturierte Flüsse von Aufmerksamkeit, Gedächtnis und Interpretation,
  \item Zyklen des Verstehens, der Vorhersage und des Lernens,
  \item eine interne kognitive Zeitskala $\tau(K_6)$.
\end{itemize}

Dieser Übergang ermöglicht das Auftreten von Bedeutung, Schlussfolgerung,
Repräsentation und Modellbildung — Fähigkeiten, die nicht auf neuronale
Dynamik reduziert werden.

% ----------------------------------------------------------------------
\subsubsection{1. Stufen und ihre Rollen}

\paragraph{K5 (neuronales Kontinuum).}
$K_5$ liefert:
\begin{itemize}
  \item spike-basierte Propagation und Synchronisation,
  \item Kanaldynamik und Achsen der Membranspannung,
  \item Erregbarkeitsschwellen $\Theta_{\mathrm{exc}}$,
  \item neuronale Flüsse $J_5$ und Propagationszyklen,
  \item Netzwerk-Konnektivität und Oszillationsmodi.
\end{itemize}

Allerdings fehlen in $K_5$:
\begin{itemize}
  \item semantische Struktur,
  \item abstrakte Repräsentationen,
  \item Schwellen auf Basis von Kohärenz oder Widerspruch,
  \item stabile konzeptionelle Attraktoren.
\end{itemize}

\paragraph{K6 (kognitives Kontinuum).}
$K_6$ führt ein:
\begin{itemize}
  \item konzeptionelle Achsen $A^{c}$, die Unterscheidungen zwischen Bedeutungen
        festlegen,
  \item Potentiale $P^{c}$ für Salienz, Überzeugungsstärke, Relevanz,
  \item Schwellen $\Theta^{c}$ für Kohärenz, Widerspruch,
        kognitive Überlastung und Zusammenbruch,
  \item Flüsse $J^{c}$: Aufmerksamkeit, Gedächtnis, Interpretation, Argumentation,
  \item Zyklen $C^{c}$: Aufmerksamkeit, Verstehen, Lernen,
  \item Maß $k(K_6)$ für kognitive Kohärenz.
\end{itemize}

Damit ist $K_6$ die erste Stufe, die ihre eigenen internen Zustände modellieren
kann.

% ----------------------------------------------------------------------
\subsubsection{2. Gemeinsame und vererbte Achsen sowie Schwellen}

\paragraph{Vererbung.}
Neuronale Synchronisationsmuster liefern das Substrat für konzeptionelle Achsen:
\[
A_{\mathrm{sync}}(K_5) \rightarrow A^{c}(K_6).
\]
Oszillatorische und Konnektivitätsmuster werden zu semantischen Unterscheidungen.

\paragraph{Neue Achsen.}
$K_6$ führt ein:
\begin{itemize}
  \item $A^{c}$ — konzeptionelle Differenzierungsachsen,
  \item $A_{\mathrm{belief}}$ — Grad von Überzeugung,
  \item $A_{\mathrm{value}}$ — Relevanz- oder Nutzengradienten,
  \item $A_{\mathrm{model}}$ — interne Modellkoordinaten.
\end{itemize}

\paragraph{Neue Schwellen.}
$\Theta^{c}$ umfasst:
\begin{itemize}
  \item \textbf{Kohärenzschwellen} — minimale Verträglichkeit von Konzepten,
  \item \textbf{Widerspruchsschwellen} — obere Grenze der Spannung innerhalb
        eines Modells,
  \item \textbf{Überlastschwellen} — Grenzen für Aufmerksamkeit oder Gedächtnisfluss,
  \item \textbf{Zusammenbruchs-Schwellen} — $\Theta_{\mathrm{collapse}}$, bei der
        $\Omega(K_6)$ leer wird, wenn Widersprüche Grenzen überschreiten.
\end{itemize}

Verletzung von $\Theta_5$ kollabiert Erregbarkeit und verhindert damit $K_6$.

% ----------------------------------------------------------------------
\subsubsection{3. Cross-Level-Flüsse, Zyklen und Spannungen}

\paragraph{Flüsse.}
Kognitive Flüsse $J^{c}$ erweitern neuronale Flüsse $J_5$:
\[
J^{c} = (J_{\mathrm{attention}},\ J_{\mathrm{memory}},\
        J_{\mathrm{interpretation}},\ J_{\mathrm{argument}}).
\]
Beispiele:
\begin{itemize}
  \item Aufmerksamkeitsverschiebungen, die neuronale Synchronisation modulieren,
  \item Gedächtniskonsolidierung, die Netzwerk-Attraktoren formt,
  \item interpretative Flüsse, die $P^{c}$ neu gewichten,
  \item argumentierende Flüsse, die konzeptionelle Achsen restrukturieren.
\end{itemize}

\paragraph{Zyklen.}
$K_6$ enthält stabile Zyklen:
\begin{align*}
C_{\mathrm{attention}} &: \text{Orientierung} \rightarrow \text{Fokus}
    \rightarrow \text{Loslösung} \rightarrow \text{Neuorientierung},\\[4pt]
C_{\mathrm{understanding}} &: \text{Vorhersage} \rightarrow \text{Input}
    \rightarrow \text{Korrektur} \rightarrow \text{aktualisiertes Modell},\\[4pt]
C_{\mathrm{learning}} &: \text{Aktivierung} \rightarrow \text{Fehler}
    \rightarrow \text{Anpassung} \rightarrow \text{Stabilisierung}.
\end{align*}

\paragraph{Spannung.}
Cross-Level-Spannung lautet:
\[
T_{5,6} = f(\Theta^{c} - P^{c},\ J^{c} - J_5,\ A^{c} - A_{\mathrm{sync}}).
\]
Überschreitet $T_{5,6}$ die dimensionsbezogene Schwelle, bricht die konzeptionelle
Kohärenz ein und das Kontinuum fällt auf ein non-semantisches neuronales Regime
zurück.

% ----------------------------------------------------------------------
\subsubsection{4. Geburts- und Todesbedingungen über die Stufen}

\paragraph{Geburt von \texorpdfstring{$K_6$}{K_6}.}
Das kognitive Kontinuum erscheint, wenn:
\[
T_5 > \Theta_{5,\mathrm{dim}}
\quad\text{und}\quad
A^{c} \in M_6 \setminus A(K_5),
\]
das Entstehen konzeptioneller Differenzierung beschreibt, die nicht auf
neurale Achsen reduzierbar ist.

Erweiterung des Zustandsraums:
\[
\Omega(K_6) = \Omega(K_5) \cup \Delta\Omega_{\mathrm{cog}}.
\]

\paragraph{Todesausbreitung.}
Wenn $\Theta_5$ versagt (Verlust der Erregbarkeit), wird $K_6$ unmöglich.

Wenn $\Theta^{c}$ versagt (Widerspruch, Überlast, Kollaps), stirbt $K_6$, während
$K_5$ bestehen bleibt:
\[
\Omega(K_6) \rightarrow \varnothing,
\qquad
\Omega(K_5) \text{ intakt}.
\]

% ----------------------------------------------------------------------
\subsubsection{5. Konzeptionelle Beispiele}

\paragraph{Bildung semantischer Achsen.}
Wiederholte Synchronisationsmuster in $K_5$ stabilisieren sich zu konzeptionellen
Unterscheidungen in $K_6$ und erzeugen Achsen wie:
\[
A^{c} = \{\text{Objekt vs. Hintergrund},\ \text{wahr vs. falsch},\
         \text{Ursache vs. Wirkung},\ \ldots\}.
\]

\paragraph{Interne Modelle.}
Prädiktive Zyklen erzeugen stabile Strukturen in $\Omega(K_6)$, die Hypothesen,
Konzepte und Erwartungen repräsentieren.

\paragraph{Kognitiver Kollaps.}
Wenn:
\[
T_{5,6} > \Theta_{\mathrm{collapse}},
\]
verliert der konzeptionelle Raum Kohärenz und $\Omega(K_6)$ schrumpft auf Null.

Daraus ergibt sich, dass der K5→K6-Übergang in komprimierter Form als Entstehen
konzeptioneller Struktur, Bedeutung und interner Modelle aus neuronaler
Dynamik und Synchronisation zusammengefasst werden kann.
